\section{Self-evaluation}\label{self-evaluation}

\subsection{Valerio Giannini}\label{self-evaluation-valerio}

\paragraph{Role.} I am the sole author and developer of KompozeR: every
bounded context described in \cref{design} --- the gateway and all eight
domain services, the Vue frontend, the Docker Compose and Kubernetes
deployments, the observability stack, and the three tiers of automated tests
(\cref{validation}) --- was designed and implemented individually, with no
division of labour to describe. The project spans roughly five months of
history (from the first architectural sketches in May to the final report
revisions in July) across some 150 commits, and grew out of an earlier
single-service coursework exercise into the full microservice system
documented here, so a meaningful part of the effort went into migrating
existing functionality onto the database-per-service, event-driven
architecture of \cref{architecture} rather than writing every feature from a
blank page.

\paragraph{Strengths.} Working alone on all layers forced the architecture
to stay coherent end-to-end: the same hexagonal layering
(\texttt{domain}/\texttt{useCases}/\texttt{adapters}) and the same
conventions (UUID identifiers, cents-based pricing, OCC via a
\texttt{version} field, \cref{implementation}) are applied identically
across all nine backend components, because there was only one person
deciding them. I am satisfied in particular with the collaborative-editing
subsystem (Lamport clocks, field-level last-writer-wins, checkpoint/replay
recovery, \cref{fault-tolerance}) and with the CAD replica set failover ---
these are the two places where a genuinely distributed-systems mechanism,
not just a description of one, is implemented and automatically tested
(\cref{validation}), including under an actual killed pod on Kubernetes. I
also kept a disciplined commit history (consistent \texttt{FEAT}/\texttt{FIX}/
\texttt{TEST}/\texttt{DOC} prefixes throughout) and a real, if manually
triggered, three-tier test suite (44 unit test files plus eight end-to-end
specs) rather than relying on manual testing alone.

\paragraph{Weaknesses.} Being both the sole developer and the sole reviewer
is itself a limitation: there was no second pair of eyes to catch design
inconsistencies before they were written down, and several of the gaps
surfaced while writing \cref{design} were only noticed at that point rather
than during implementation --- most notably the absence of any continuous
integration pipeline (tests are run manually, \cref{validation}), the
complete lack of an automated test suite on the frontend, the single-attempt
(no-retry) behaviour of most inter-service REST calls, and a few rough edges
in the deployment configuration (a stale \texttt{.env.example} that no
longer matches the services actually running, and every Kubernetes
Deployment fixed at one replica even though the stateless tier is designed
to scale horizontally). Secret management for local/demo environments is
also weaker than it should be and is the first thing I would fix outside the
scope of this report. Time-boxing the project around the course deadline
also meant prioritising breadth --- getting all eight bounded contexts and
both real-time subsystems working end-to-end --- over hardening any single
one of them further; \cref{future-works} lists the concrete consequences of
that trade-off.
